\section{書かれた一粒子作用}
\label{sec:written}

地位は\dstrec{}である。ソースは \texttt{DualRotorDynamics} と \texttt{DualRotorFlow} である。実装は数値 ODE を先に書かない。この章の関数は閉形式を返す。

\subsection{密度とオイラー--ラグランジュ方程式}

Lean が定義している密度は一軸のものである。
\[
L=\frac12\dot\phi^2-\frac12\dot\theta^2-\frac m2(\phi-\theta)^2
\]
識別子は \texttt{paperLagrangian} である。この密度に軸をまたぐ項は無い。三軸を同時に進めるときは、同じ密度を軸ごとの初期値で別々に評価する。

正しいオイラー--ラグランジュ方程式 \texttt{PaperActualEL} は
\[
\ddot\phi=-m(\phi-\theta),\qquad
\ddot\theta=-m(\phi-\theta)
\]
である。したがって
\[
\ddot\delta=0,\qquad
\ddot\sigma=-2m\delta.
\]
書かれた方程式は \(\ddot\delta+m\delta=0\) ではない。定理 \texttt{paperActualEL\_not\_oscillator} の証人は
\[
(m,\phi,\theta,\ddot\phi,\ddot\theta)=(1,1,0,-1,-1)
\]
であり、この点では \(\ddot\delta=0\) かつ \(-m\delta=-1\) である。第\ref{sec:working}章の同符号模型はさらに別の方程式 \(\ddot\delta+2m\delta=0\) であり、ここへ混ぜない。

\subsection{閉形式}

初期値を \(\delta(0)=\delta_0\)、\(\dot\delta=\nu\)（定数）、\(\sigma(0)=\sigma_0\)、\(\dot\sigma(0)=\dot\sigma_0\) とする。解は
\begin{align}
\delta(t)&=\delta_0+\nu t,\\
\dot\sigma(t)&=\dot\sigma_0-2m\delta_0 t-m\nu t^2,\\
\sigma(t)&=\sigma_0+\dot\sigma_0 t-m\delta_0 t^2-\frac{m\nu}{3}t^3
\end{align}
である。識別子は \texttt{writtenDelta}、\texttt{cubicRapiditySpeed}、\texttt{writtenSigma} である。角へ戻すときは
\begin{verbatim}
phi(t) = (sigma(t) + delta(t)) / 2
theta(t) = (sigma(t) - delta(t)) / 2
\end{verbatim}
とする。\(\ddot\phi=\ddot\theta=-m\delta(t)\) が各時刻で成り立つことを、実装後の検査に入れる。

\subsection{運動量とヤコビ積分}

共役運動量に剛性 \(m\) は掛からない。
\begin{align}
p_\sigma&=\frac12\dot\delta=\frac12(\dot\phi-\dot\theta),\\
p_\delta&=\frac12\dot\sigma=\frac12(\dot\phi+\dot\theta).
\end{align}
識別子は \texttt{paperMomentumSigma} と \texttt{paperMomentumDelta} である。書かれた流れの上では \(p_\sigma\) は保存し、\(\dot p_\delta=-m\delta\) である。

ヤコビ積分は
\[
E=\frac12\dot\phi^2-\frac12\dot\theta^2+\frac m2\delta^2
=\frac12\dot\sigma\dot\delta+\frac m2\delta^2
\]
であり、識別子は \texttt{paperEnergy} である。\texttt{PaperActualEL} の上で \(\dot E=0\) である。閉形式の上では、すべての時刻で
\[
E(t)=\frac12\nu\dot\sigma_0+\frac m2\delta_0^2
\]
である。定理は \texttt{written\_flow\_energy} である。\(E\) は不定であり、正にも負にもなりうる。この保存は、振動数の存在を意味しない。

\subsection{呼び出し境界}

\begin{verbatim}
function written_at(seed, t) -> (phi, theta, phidot, thetadot):
    delta = seed.delta0 + seed.nu * t
    sigmadot = seed.sigmadot0
             - 2 * seed.stiffness * seed.delta0 * t
             - seed.stiffness * seed.nu * t * t
    sigma = seed.sigma0
          + seed.sigmadot0 * t
          - seed.stiffness * seed.delta0 * t * t
          - (seed.stiffness * seed.nu / 3) * t * t * t
    phi = (sigma + delta) / 2
    theta = (sigma - delta) / 2
    phidot = (sigmadot + seed.nu) / 2
    thetadot = (sigmadot - seed.nu) / 2
    return (phi, theta, phidot, thetadot)
\end{verbatim}

この関数は力も位置も受け取らない。駆動 \(u\)、振動子、暴走、任意ポテンシャルは第\ref{sec:working}章の別モジュールである。
